\documentclass[12pt]{article}
\usepackage[utf8]{inputenc}
\usepackage[T1]{fontenc}
\usepackage{hyperref}
\title{Fipa Language Resources for K-12 Curriculum Development}
\author{Compiled by Claude Sonnet 5.5 (claude-sonnet-5-5) for Andr\'as Kornai}
\date{}
\begin{document}
\maketitle

\section{Introduction}
Fipa (fip) is a Bantu language of the Rukwa Region of southwestern Tanzania, with about 200,000 speakers. It is not known to be taught in schools. Very little was found; no dictionary, primer or textbook appeared in the search results.

\section{Primary Educational Resources}
None found.

\section{Linguistic and Reference Materials}
\begin{itemize}
\item Overview of the Fipa language: \url{https://en.wikipedia.org/wiki/Fipa_language}
\item Colloquium abstract, ``Fipa: a first look at dialectal variation and syntactic variation'' (Humboldt-Universit\"at zu Berlin): \url{https://iaaw.hu-berlin.de/de/afrika/linguistik-und-sprachen/veranstaltungen/afrikalinguistischeskolloquium/papers-sommersemester-2010/fipa-a-first-look-at-dialectal-variation-and-syntactic-variation}
\end{itemize}

\section{Religious and Community Resources}
\begin{itemize}
\item Joshua Project people-group page, reporting that a New Testament exists but not a full Bible: \url{https://joshuaproject.net/people_groups/11211/TZ}
\end{itemize}

\section{Digital and Audio Resources}
\begin{itemize}
\item Scripture Earth listing for Fipa (New Testament and audio recordings): \url{https://www.scriptureearth.org/00eng.php?iso=fip}
\end{itemize}

\section{Curriculum Development Context}
Fipa-medium school materials were not found. A curriculum would have to be built from scratch, starting from the New Testament text and linguistic descriptions.

\end{document}